\section{n3. Irreducible self-crossing has no greatest C/D minorant; exact canonical error is ceil(3m/2)/m-1.}
\textit{Supports: G331, n1, n2.}

An irreducible self-crossing example, with exact canonical ideals in EVERY degree. Use ENTIRE qualified G331 and n1--n2. Let k be ANY field with char(k)!=2 and d\textgreater{}=2, let z=(z\_1,...,z\_(d-1)), and
R=k[u,v,z]\_(u,v,z)/(v\textasciicircum{}2-u\textasciicircum{}2(1+u)), m=(u,v,z).
This is an excellent DOMAIN of dimension d. The element 1+u is not a square in k(u,z), since its prime divisor has odd order. Thus the displayed quadratic is irreducible; its parametrization u=t\textasciicircum{}2-1, v=t(t\textasciicircum{}2-1) embeds R in its fraction field. Write S\_0=k[t,z]. It is normal, finite over the unlocalized R\_0 via integral t with t\textasciicircum{}2=1+u, and has the same fraction field since t=v/u. Hence it is the normalization. Localizing at R\_0\textbackslash{}m gives a NORMAL DOMAIN S with exactly two closed centres M\_+=(t-1,z), M\_-=(t+1,z). It is the semilocalization of k[t,z] at these centres: every element outside their union is a unit in the finite algebra over local R, proving the identification with that semilocalization. S is not a product of two rings.

All conductor claims are exact. Put A=k[u,z]. Then S\_0=A direct-sum tA as A-module and R\_0=A direct-sum tuA. For a+t b in S\_0, membership in R\_0 forces b in uA; multiplication by t then forces a in uA as well. Conversely uS\_0 multiplies S\_0 into R\_0. Thus the conductor is uS\_0 before localization and uS afterwards, since the conductor is the annihilator of the finite module S\_0/R\_0 and annihilators of finite modules localize. Its target ideal is (u,v). The quotient is R/c=k[z]\_(z), while S/c is TWO copies of that ring, because t\textasciicircum{}2-1=(t-1)(t+1) with coprime factors. Exactly R=k[z]\_(z)+uS: subtract the lifted quotient residue, leaving uS. This is a conductor fibre product, not a decomposition of the normalization domain. Moreover mS=(u,z); on either normal local centre this is its maximal ideal since the opposite t-sign factor is a unit.

Consequently every normalized target semivaluation has its source lift(s) on these two centres. The finite-u stratum has unique lifts. At infinite u-value, the ENTIRE quotient-boundary semivaluation on k[z] has two lifts, one at each centre, and these are exactly the identifications by G331. For d=2 it is one conductor point o with z\_1=1; for d\textgreater{}2 it retains every support of the z-quotient. This proves the marked-centre hypothesis of n2 in an irreducible nonnormal ring.

There is a fully explicit obstacle requiring no unknown Izumi constant. Let ell=v-u=u(t-1) in R, I=(ell)+m\textasciicircum{}2, and F=(1/2)v(I)-2v(m). This is a continuous signed ACTUAL primary-model evaluation. On the plus centre, put s=w(t-1), Z=min\_j w(z\_j), so min(s,Z)=1 and w(u)=s; t+1 is a unit. Then w(ell)=2s\textgreater{}=2, hence F=-1 on its ENTIRE boundary. On the minus centre put r=w(t+1), with min(r,Z)=1, and t-1 a unit. There w(ell)=r and F=(1/2)min(r,2)-2, between -3/2 and -1. All conductor values are -1. These identities hold when coordinates are infinite; ideal I is proper primary and its evaluations are finite.

The generic ordinary-order valuation nu\_- on the regular local minus centre has r=Z=1 and F(nu\_-)=-3/2. Every normalized semivaluation on that centre dominates ordinary order: all regular parameters have value\textgreater{}=1, every element of the e-th maximal power has value\textgreater{}=e, and units have value0. Thus for every unit-or-primary target ideal w(J)\textgreater{}=nu\_-(J). Negative REAL combinations, G331 uniform comparison and ACTUAL decreasing infima give g(rho w)\textless{}=g(rho nu\_-) for EVERY target geometric C/D function on the entire minus image, including all conductor supports. Every minorant of F is therefore \textless{}=-3/2 there.

For n\textgreater{}=1 define the ACTUAL conductor element j\_n=u(t+1)\textasciicircum{}(n-1) in R and the proper primary ideal T\_n=m\textasciicircum{}(3n)+(j\_n)m\textasciicircum{}(2n), with q\_n=-v(T\_n)/(2n). On the plus centre, w(j\_n)=s, so q\_n=max(-3/2,-1-s/(2n)). On the minus centre, w(j\_n)=nr\textgreater{}=n, so q\_n=-3/2. They are \textless{}=F on every source and target support. The negative effective primary divisors -E\_(T\_n)/(2n) are relatively ample and positive on EVERY contracted integral curve before any limit. The constant -3/2 is also a competitor. Hence the COMPLETE raw C/D supremum is
p=-1 on the plus-centre finite-u locus, and p=-3/2 on the entire minus image, including every conductor point.
No generic-support restriction is present in this formula.

At the plus centre the monomial valuations w\_R(t-1)=R, w\_R(z\_j)=1, R-\textgreater{}infinity, converge pointwise on every element of its polynomial local ring to the valuation killing t-1 and having ordinary z-order. Every fixed unit denominator has value0. Their target restrictions converge to a conductor point, where the raw value drops by 1/2. Thus p is not USC and neither a greatest C nor a greatest D minorant exists. This is NONNORMAL failure on an irreducible germ, not normal CE failure. Its USC regularization has value -1 at that conductor point and -3/2 at nu\_-: the open test v(ell)\textless{}2 around nu\_- meets only the minus image. It violates g(o)\textless{}=g(nu\_-), so regularization does not restore geometric positivity.

The canonical filtration has an exact finite algebraic description. For EVERY integer m\textgreater{}=1 put c\_m=ceil(3m/2), and in S put
J\_m=(t-1,z)\textasciicircum{}(m-1)(t+1,z)\textasciicircum{}(c\_m-1).
The two semilocal ideals are comaximal since 2 is invertible; their product equals their intersection. Then
A\_m(F)=(z)\textasciicircum{}(c\_m) R + u J\_m,
where uJ\_m is contained in R by the conductor. This is a genuine R-ideal with explicit finite generators: all z-monomials of total degree c\_m, and u t\textasciicircum{}epsilon times every product of a degree-(m-1) monomial in (t-1,z) and a degree-(c\_m-1) monomial in (t+1,z), for epsilon=0,1. S=R+Rt, so these generators span exactly uJ\_m over R. Each polynomial is an actual downstairs element; use t\textasciicircum{}2=1+u and ut=v to write it in u,v,z.

Proof of the identity: testing ordinary order at the plus and minus centres forces respectively order\textgreater{}=m and order\textgreater{}=c\_m. Those conditions suffice on EVERY source support because normalized values dominate order and F is respectively -1 and \textgreater{}=-3/2. Every f in R has a decomposition f=g(z)+u a, with g in the local z-ring and a in S. The minus order condition forces g in (z)\textasciicircum{}(c\_m), by reducing modulo u; the quotient regular local ring has its usual power filtration. Subtract this residue lift. Since u is a regular parameter on BOTH normal centres, the remaining u a has the required orders exactly when a is in the localized (m-1)-st plus power and (c\_m-1)-st minus power. In a semilocal ring these local conditions are their intersection, equal to J\_m. Conversely the displayed ideal satisfies both conditions. This argument includes localized unit denominators and all cancellation; it is not an inferred monomial filtration.

Upon localization at the plus centre, the other centre's power is the unit ideal. Thus on EVERY plus support
h\_m=-min(c\_m Z,s+m-1)/m.
At every minus support, including conductor supports, h\_m=-c\_m/m. Every h\_m has the ACTUAL negative effective ample primary presentation from n1. These formulas prove pointwise convergence in EVERY degree to p. The EXACT uniform error is
||p-h\_m||=c\_m/m-1 \textgreater{}=1/2.
Indeed on the plus finite-u locus the gap is at most c\_m/m-1, attained by Z=1 and finite s\textgreater{}=c\_m-m+1. On the minus image it is c\_m/m-3/2, which is 0 in even degree and 1/(2m) in odd degree. Therefore the largest gap is 1/2 in even degree and 1/2+1/(2m) in odd degree. The canonical algebra is not finitely generated by n1's proved forward implication. Failure persists for uniform perturbations \textless{}1/4 by q\_n-epsilon and the same ordinary-order inequality.

The same construction works in every dimension d\textgreater{}=2 over every field of characteristic different from2, as proved above. It also works for the genuinely mixed-characteristic irreducible SURFACE R=V[u,v]\_(pi,u,v)/(v\textasciicircum{}2-u\textasciicircum{}2(1+u)), for any excellent DVR V with 2 invertible. Its normalization is the semilocal domain V[t] at (pi,t-1),(pi,t+1), the conductor is uS, the quotients are V and V times V, and mS=(pi,u). The same free rank-two module calculation proves normalization and conductor; irreducibility follows over Frac(V)(u), and regularity/normality of V[t] supplies the source. Replace z by pi throughout. Ordinary (pi,t+1)-order is a valuation since its associated graded ring is a polynomial domain, and every normalized support dominates it. The same explicit F,T\_n,A\_m,h\_m,p and exact gap formulas hold. No claim is made when2 is not invertible. The characteristic-two displayed field equation degenerates and is not an automatic transfer of this example.

All expressions are homogeneous on rays. Under a transported primary obstacle the canonical ideals are unchanged and profiles and p divide by the bounded positive gauge factor; greatest-minorant failure is gauge invariant. Completion transport holds, although the irreducible nonnormal ring need not have a domain completion. No analytic irreducibility assertion is made for R. Literal radial comparison, products, energy and full operators remain separate. Artifacts n3.md, n3.tex and reports/diagnostics.json; rerun python diagnostics.py; python assemble.py; python verify\_delivery.py. Finite checks below are observations on explicit input domains; the universal proofs are the argument here.

\[
 R=k[u,v,\mathbf z]_0/(v^2-u^2(1+u)),\quad u=t^2-1,\quad v=tu.
\]
\[
 F=\tfrac12v((v-u)+\mathfrak m^2)-2v(\mathfrak m),\qquad
 c_m=\lceil3m/2\rceil.
\]
\[
 A_m=(\mathbf z)^{c_m}R+u(t-1,\mathbf z)^{m-1}(t+1,\mathbf z)^{c_m-1},
 \qquad\|p-h_m\|_\infty=c_m/m-1\ge\tfrac12.
\]

